\[ \Th \nexists \Lambda\in\mathbb{R} \mid \Lambda=\sup\mathbb{N} \]

\begin{enumerate}[label=(\roman*)]
    \item Supponiamo per assurdo che 
        \[
            \exists \Lambda\in\mathbb{R} \mid \Lambda=\sup\mathbb{N}
        \]
    \item Essendo $\Lambda$ il minimo dei maggioranti (poiché estremo superiore),
    allora $\Lambda - 1$ non può essere un maggiorante per $\mathbb{N}$, ovvero
        \[
            \exists n\in\mathbb{N}\mid\Lambda-1<n,
            \exists n\in\mathbb{N}\mid\Lambda<n+1.
        \]

    \item Ma sapendo per \defref{def:naturals} che $n\in\mathbb{N}\Rightarrow n+1\in\mathbb{N}$,
    allora
        \[
            \exists n\in\mathbb{N}\mid\Lambda<n
        \]
    \item Dunque $\Lambda=\sup\mathbb{N}$ non può essere un maggiorante per $\mathbb{N}$, non può quindi
    nemmeno essere l'estremo superiore.

    \item La negazione della tesi porta a una contraddizione.
    \item La tesi $\nexists \Lambda\in\mathbb{R} \mid \Lambda=\sup\mathbb{N}$ deve essere
    necessariamente vera.
    \item Quindi $\mathbb{N}$ non è superiormente limitato.
    \item Dunque $\sup\mathbb{N} = +\infty$.
\end{enumerate}
\qed